\section{Discussion}
\label{sec:discussion}

\paragraph{Hypotheses.}
\begin{itemize}
\item \textbf{H1, holds.} Muting or weighting microphones gains under
  1\,dB (\Gainmute{} and \Gainmutesoft\,dB, WER \DiffMuteVsRef{} points,
  $p=$\,\PMuteVsRef). Phase-aware front ends gain 3--4\,dB.
  \begin{itemize}
  \item We predicted \emph{at least 5\,dB} for MVDR. Only the oracle
    reaches that, so the bottleneck is the noise model from 1.5\,s of
    silence, not the array.
  \end{itemize}
\item \textbf{H2, half right.} The spatial front end matters most for
  speech-like interferers, as predicted. But DeepFilterNet3 alone does not
  ``handle'' stationary noise for the recognizer either: it improves ESTOI
  and leaves the WER flat or worse.
\item \textbf{H3, holds, but the optimum does not move.} WER is minimal
  from ``off'' to 12\,dB and grows beyond, both with and without a
  beamformer in front.
\item \textbf{H4, holds.} The robot's independent 64\,ms chunks cost 6--11
  points.
\end{itemize}

\paragraph{Falsification test.} We set this test in advance:
MVDR$\to$DeepFilterNet3 against DeepFilterNet3 alone at noise angles
$\geq$\,90$^\circ$. On those 44 scenes the difference is $-26.3$ WER
points, 38\,\% relative, far beyond the 10\,\% threshold. The difference is
the same below 90$^\circ$, and over all 0\,dB scenes it is
\DiffRtfDfnVsRefDfn{} points ($p=$\,\PRtfDfnVsRefDfn). The 4-microphone front end is worth adding.

\paragraph{Recommendations for the robot.}
\begin{enumerate}
\item Read the raw channels, not ch0.
\item Run an RTF-steered MVDR whose noise covariance is updated in the
  silence before each utterance. The voice activity detector already
  provides that silence.
\item Run DeepFilterNet3 continuously, or not at all, at 12\,dB or below.
  Never run it without a limit in front of the recognizer.
\end{enumerate}
The setting \texttt{atten\_lim\_db}\,=\,0 deserves a warning, because the
library reads it as ``no limit''.

\paragraph{Limitations.}
\begin{itemize}
\item One talker, one room and one session. The talker's angle marks were
  approximate, so separations come from the clean takes.
\item Scenes are mixed. The live check shows they are optimistic by 1--2\,dB
  of SNR gain, but the two live takes agree with them in direction and size.
\item Whisper \texttt{small} with greedy decoding; a larger model is more
  noise-robust and may shrink every gap.
\item The 0\,dB scenes are hard on purpose: the absolute WERs are high, the
  differences are what matters.
\item The cafeteria live take could not be checked for superposition: a
  bug in the noise generator replayed a different excerpt.
\end{itemize}

\paragraph{Project~15.} The same session gives project~15 its missing
hardware numbers:
\begin{itemize}
\item \textbf{Orientation.} The array is mirrored and rotated by
  \PFifteenOffset$^\circ$ relative to the talker's marks.
\item \textbf{Direction.} The leave-one-out median error is
  \PFifteenLooErr$^\circ$, bounded by the marks' accuracy.
\item \textbf{Distance fails.} No distance cue correlates with distance
  (|Spearman|\,$\leq$\,\PFifteenSpearmanMax). The best model's error is
  \PFifteenMaeBest\,m, against \PFifteenMaeMean\,m for predicting the mean.
\end{itemize}
